\section{Exact interfaces with the area-law companion}
\label{app:companion}

The companion \cite{RepoAreaLaw}, \emph{A two-dimensional area law from a
global spectral gap}, is used on the original ground state $\Omega$.
This appendix records the hypotheses and the division of work between the
two articles. A subregion is always a subsystem of that state; it is never
assigned a new gap assumption.

\paragraph{Hamiltonians, metrics, and entropy.}
The companion permits any finite induced domain $\Gamma\subset\mathbb Z^2$
with local dimension $q$ and one Hermitian term $h_X$ of norm at most $J$
for each nonempty support of induced-graph diameter at most a fixed $R$.
It assumes a unique unit ground vector and
$H-E_0I\ge\Delta(I-|\Omega\rangle\langle\Omega|)$.
Constants depend on $q,R,J,\Delta$ and explicitly fixed auxiliary
parameters, uniformly in the domain and Hamiltonian. All entropy logarithms
are natural. Its Theorem~1.1 and Corollary~1.2 give the area inequality
stated in Proposition~\ref{prop:area-law}. Here $\Gamma=\Lambda_L$ and
$R=1$, exactly covering \eqref{eq:model}.

\paragraph{Positive replacement.}
The Proposition ``Positive replacement of the Hamiltonian'' in
\cite[Section~4]{RepoAreaLaw} gives, for every fixed integer $k\ge1$ and
$\alpha=k/(k+1)$, positive contractions $k_i$ at the original interaction
anchors, of uniformly bounded anchor multiplicity, satisfying
\[
 k_i\Omega=0,\qquad
 \sum_i k_i\ge c_*^{-1}(H-E_0I)\ge(\Delta/c_*)(I-|\Omega\rangle\langle\Omega|),
 \qquad
 \|k_i-\mathbb E_{N_r(a_i)}k_i\|\le C e^{-cr^\alpha}.
\]
Here $\mathbb E_{N_r(a_i)}$ is conditional expectation onto the induced-graph
ball, for every integer $r\ge0$; $c_*,C,c$ may depend on $k$ and the fixed physical parameters. This
input underlies the companion amplification. It does not replace the
original Hamiltonian by a finite-range frustration-free Hamiltonian.

\paragraph{Templates, collars, and amplification.}
Definition ``A template'' and Proposition ``A sublinear collar bound'' in
\cite[Section~9]{RepoAreaLaw} use a cut $A\subseteq\Gamma$, its crossing-edge
endpoints $Z_A$, and a nonempty ambient lattice set
\[
 T=\bigcup_{a=1}^{N_0}(P_a\cap\mathbb Z^2),\qquad
 n\ge C_{\rm tpl}N_0(s_0+1),\quad \diam_\infty T\le n,
 \quad \dist_\infty(T,Z_A)>4D_0s_0.
\]
The $P_a$ are closed convex rectangles or triangles of sup diameter at most
$s_0$, with side slopes $0,\infty,1,-1$; $n,s_0$ are positive integers and
$s_0\ge1$. With $\epsilon=10^{-5}$, set
$\ell=\lfloor n^{1-\epsilon}\rfloor\le s_0$,
$T_\ell=T+[-\ell,\ell]^2_{\mathbb Z}$, $X=A\cap T$,
$S_0=A\cap T_\ell$, and $Y=S_0\setminus X$.
At all sufficiently large scales the input guarantees
\[
 I_\Omega(X:\Gamma\setminus S_0)\le Cn^{1-\epsilon},\qquad
 |S_0|\le Cn^2,\qquad
 |\partial_\Gamma X|,|\partial_\Gamma Y|,|\partial_\Gamma S_0|\le Cn.
\]

The Proposition ``Amplification'' in \cite[Section~10]{RepoAreaLaw}
requires these three bounds with any fixed common constant $C_0$ and
$\alpha=k/(k+1)>1-10^{-6}$. Put $\beta_0=1-2\cdot10^{-6}$.
For every fixed integer $M\ge1$ it supplies a radius
\[
 r\le C_M n^{(1-\epsilon)/\alpha},\qquad
 \ell+2r=o(n^{\beta_0}),\qquad
 I_\Omega(X':E)\le C_M n^{-M}
\]
simultaneously for every $X'\subseteq X$ and
$E\subseteq\Gamma\setminus N_{2r}(S_0)$. The radius and constants are
independent of the remote set and total volume. If also
$\dist_\infty(T,Z_A)>\ell+n^{\beta_0}$, then
$N_{2r}(S_0)\subseteq A$. In Proposition~\ref{prop:coarse} we take
$A=\Lambda$, so $Z_A=\varnothing$, choose $s_0,n$ comparable to the target
scale with fixed ratios, and use $M=1$. A graph neighborhood lies in the
corresponding ambient enlargement; this is the precise conversion used to
obtain the ambient-separation statement.

\paragraph{The two-family interface and its adaptation.}
The Lemma ``Two-family entropy cancellation'' in
\cite[Section~11]{RepoAreaLaw} states that, for a pure state and a partition
$A=D\mathbin{\dot\cup}\bigcup_iX_i$ into two ordered families, bounds
\[
 I\left(X_i:A^c\cup\bigcup_{j<i,\,f_j=f_i}X_j\right)\le\delta_i
 \quad(f_i\in\{0,1\})
 \quad\Longrightarrow\quad
 S(A)\le S(D)+\tfrac12\sum_i\delta_i.
\]
Its Proposition ``Two families with separated birth templates'' constructs
such regions for every cut with $b=|\partial_\Gamma A|>0$ and any fixed lower
scale $n_*$. It gives $|D|\le Cb$, $X_i\subseteq A\cap T_i$, admissible
separated templates with $n_i\ge n_*$ and
$\lfloor n_i^{1-\epsilon}\rfloor\le s_i$, and
\[
 \dist_\infty(T_i,Z_A)>2n_i^{\beta_0},\qquad
 \dist_\infty\left(T_i,\bigcup_{j<i,\,f_j=f_i}X_j\right)>2n_i^{\beta_0},
 \qquad \sum_i n_i^{-100}\le Cb.
\]
The constant may depend on $n_*$ and the template constants but not on the
cut or domain. These statements are used within the companion's area-law
proof. Lemma~\ref{lem:cell-cmi} adapts their ordering principle to a cell
with arbitrary exterior conditioning and proves that local variant in
full. In particular, we do not infer a conditional-information bound merely
by conditioning an ordinary mutual-information inequality.

The companion's replica limit is taken at each fixed finite physical
instance and schedule, before the geometric scale tends to infinity.
Its exported bounds concern the original one-copy state with the uniform
constants above. The later tapering construction, adaptive patch
constraints, density compression, encoded protocol, and routing are proved
in this article.
